% !TEX program = xelatex
% 数学物理方法2 · §7.2 · 课件 p19 作业第 4 题 · 解答卡（2026-11-09）
\documentclass[10pt]{ctexart}
%=====================================================================
% theme-violet.tex —— 讲义精讲 · 紫花主题
% 用法：\documentclass[11pt]{ctexart}
%       \input{assets/theme-violet}
%       \begin{document} ... \end{document}
% 只用 ctexart 自带能力 + xcolor/enumitem/fancyhdr/pifont/amsmath，
% 不依赖 tikz/tcolorbox，保证 MiKTeX 可直接编译。
%=====================================================================

\usepackage[a4paper,left=2.0cm,right=2.0cm,top=2.1cm,bottom=2.3cm,%
            headsep=12pt,footskip=22pt]{geometry}
\usepackage{amsmath,amssymb}
\usepackage{needspace}
\usepackage{xcolor}
\usepackage{array}
\usepackage{enumitem}
\usepackage{pifont}
\usepackage{fancyhdr}
\usepackage[hidelinks]{hyperref}

%---------------------------------------------------------------- 配色
\definecolor{violet}{RGB}{107,78,143}      % 主色：紫
\definecolor{violetdeep}{RGB}{63,44,92}    % 深紫：标题/强调
\definecolor{petal}{RGB}{183,159,208}      % 花瓣：装饰线、花饰
\definecolor{bandviolet}{RGB}{247,243,251} % 浅紫底：档案/提示框
\definecolor{bandcream}{RGB}{252,249,242}  % 米白底：易错/注意框
\definecolor{rulepurple}{RGB}{217,207,230} % 分隔线
\definecolor{inkgray}{RGB}{92,96,110}      % 次要文字

%---------------------------------------------------------------- 字体
% 有 Noto CJK 用 Noto，没有就退回 ctex 默认（Windows 下为宋体/黑体）
\IfFontExistsTF{Noto Serif CJK SC}{\setCJKmainfont{Noto Serif CJK SC}}{}
\IfFontExistsTF{Noto Sans CJK SC}{\setCJKsansfont{Noto Sans CJK SC}}{}

%---------------------------------------------------------------- 版式
\setlist{nosep,leftmargin=1.9em}
\linespread{1.28}
\setlength{\parskip}{3.5pt}
\setlength{\parindent}{0pt}
\setlength{\tabcolsep}{5pt}
\renewcommand{\arraystretch}{1.28}

%---------------------------------------------------------------- 花饰
% \ding{95} = 五瓣小花，\ding{96} = 星形花；都来自 pifont，必定存在
\newcommand{\orn}{\textcolor{petal}{\ding{95}}}          % 行内小标
\newcommand{\ornstar}{\textcolor{petal}{\ding{96}}}      % 星形花
\newcommand{\flowerrule}{%                                % 居中花饰分隔线
  \par\vspace{7pt}%
  \noindent\makebox[\linewidth]{\color{petal}%
    \rule{0.20\linewidth}{0.6pt}\hspace{8pt}{\small\ding{95}}\hspace{4pt}%
    \raisebox{1.1pt}{\tiny\ding{95}}\hspace{4pt}{\small\ding{95}}\hspace{8pt}%
    \rule{0.20\linewidth}{0.6pt}}%
  \par\vspace{7pt}}

%---------------------------------------------------------------- 标题与节头
\newcommand{\doctitle}[3]{%   #1 主标题  #2 副标题  #3 元信息（可用 \\ 分行）
  \begin{center}
    {\color{petal}\small \orn\enspace #2\enspace \ornstar}\\[7pt]
    {\LARGE\bfseries\color{violetdeep} #1}\\[10pt]
    {\color{petal}\rule{0.55\linewidth}{1pt}}\\[4pt]
    {\small\color{inkgray} #3}
  \end{center}\vspace{2pt}}

\newcommand{\hdr}[2]{%   #1 章节号（如 §0）  #2 标题
  \par\vspace{12pt}%
  \noindent\colorbox{violetdeep}{\color{white}\bfseries\small\;#1\;}\hspace{7pt}%
  {\color{violetdeep}\large\bfseries #2}\par\nobreak
  \vspace{3pt}{\color{petal}\hrule height 1.2pt}\vspace{6pt}}

\newcommand{\subhdr}[1]{%
  \par\vspace{6pt}%
  \noindent{\color{violet}\bfseries #1}\par\vspace{1pt}}

%---------------------------------------------------------------- 提示框
\newcommand{\callout}[1]{% 浅紫底：讲义档案、读法说明
  \par\medskip\noindent
  \fcolorbox{rulepurple}{bandviolet}{\parbox{\dimexpr\linewidth-2\fboxsep-2\fboxrule}{\vspace{2pt}#1\vspace{2pt}}}%
  \par\medskip}

\newcommand{\warnbox}[1]{% 米白底：易错、边界条件
  \par\medskip\noindent
  \fcolorbox{rulepurple}{bandcream}{\parbox{\dimexpr\linewidth-2\fboxsep-2\fboxrule}{\vspace{2pt}#1\vspace{2pt}}}%
  \par\medskip}

\newcommand{\gapitem}[2]{% #1 缺口类型（紫） #2 内容
  \par\vspace{2pt}\noindent
  \hangindent=1.4em\hangafter=1
  {\color{petal}\small\ding{96}}\enspace{\bfseries\color{violetdeep}#1}\enspace #2\par}

% 题干块（\qstem）：每道例题／作业的「详解」开头必须先复述题干与出处。
% 起因：2026-10-25 第五章讲义把七道作业的题面集中放在一页、详解另起一页，
% 读者翻到「第 4 题」时看不到条件，等于一道没有题干的题目。
% 规则：详解块内自带题干，读者不翻页就能读懂题在问什么。
\newcommand{\qstem}[1]{%
  \par\medskip\noindent
  \fcolorbox{violet}{bandviolet}{%
    \parbox{\dimexpr\linewidth-2\fboxsep-2\fboxrule}{%
      \vspace{2pt}%
      {\small\bfseries\color{violet}题干}\hspace{0.8em}#1%
      \vspace{2pt}}}%
  \par\medskip}

% v2: 保护题干与开头，长题可跨页，续页补上下文。
\newcommand{\qbreak}{\needspace{6\baselineskip}}
% 独立练习：ID、来源、题干、考点、答案标签（通常 ans:ID）。
\newcommand{\practice}[5]{%
  \qbreak\subhdr{练习 #1}\qstem{（#2）#3}%
  \textbf{考点}：#4\par
  详解见附录第~\pageref{#5}~页。\par\medskip}
% 附录开头：ID、来源、条件摘要；自动产生 ans:ID 标签。
\newcommand{\answerstart}[3]{%
  \qbreak\subhdr{#1\quad 详解}\label{ans:#1}\qstem{（#2）#3}}
\newcommand{\answercontinue}[2]{%
  \needspace{5\baselineskip}\subhdr{#1\quad 续解}%
  {\small\color{inkgray}本页接续：#2}\par\medskip}
% 延迟提供默认值，让旧讲义可在导言区保留自己的 \newcommand{\step}。
\AtBeginDocument{\providecommand{\step}[1]{\needspace{3\baselineskip}\par\medskip\textbf{#1}\par}}

% 前提旁注：只放在「需要理解的那一步」后面，给陈述 + 本题代入，不重证。
\newcommand{\aside}[1]{%
  \par\vspace{1pt}\noindent
  \hangindent=1.6em\hangafter=1
  {\small\color{inkgray}%
    \hspace*{0.6em}{\color{petal}$\triangleright$}\hspace{2pt}#1}%
  \par\vspace{2pt}}

%---------------------------------------------------------------- 表格
% 本机 MiKTeX 的 colortbl 与 array 版本不匹配（\rowcolor 会报错），
% 故表格统一用「上粗紫线 + 表头 + 中细线」的三线式，不依赖 colortbl。
\newcommand{\tabtop}{\par\vspace{3pt}{\color{violet}\hrule height 1.2pt}\vspace{1pt}}
\newcommand{\tabmid}{\par\vspace{1pt}{\color{petal}\hrule height 0.5pt}\vspace{1pt}}
\newcommand{\tabbot}{\par\vspace{3pt}{\color{rulepurple}\hrule height 0.5pt}\vspace{8pt}}
\newcommand{\thead}[1]{{\bfseries\color{violetdeep}#1}}

%---------------------------------------------------------------- 页眉页脚
\pagestyle{fancy}
\fancyhf{}
\renewcommand{\headrulewidth}{0.4pt}
\renewcommand{\headrule}{\hbox to\headwidth{\color{rulepurple}\leaders\hrule height \headrulewidth\hfill}}
\fancyhead[L]{\small\color{inkgray}\songti \doctag}
\fancyhead[R]{\small\color{inkgray}\songti \coursetag}
\fancyfoot[C]{\small\color{inkgray}\orn\enspace{\color{violet}\thepage}\enspace\orn}

% 每篇讲义在导言区设这两个值
\newcommand{\doctag}{}
\newcommand{\coursetag}{}
\newcommand{\settags}[2]{\renewcommand{\doctag}{#1}\renewcommand{\coursetag}{#2}}

%---------------------------------------------------------------- 数学记号
\newcommand{\vect}[1]{\boldsymbol{#1}}
\newcommand{\dd}{\,\mathrm{d}}
\newcommand{\ddx}{\frac{\mathrm{d}}{\mathrm{d}x}}
\addtolength{\jot}{2pt}



\settags{§7.2 作业 · 热流边界条件}{数学物理方法2}

\geometry{left=1.9cm,right=1.9cm,top=1.9cm,bottom=2.0cm,headsep=10pt,footskip=18pt}
\linespread{1.24}
\setlength{\headheight}{22pt}
\setlength{\parskip}{0pt}
\setlength{\abovedisplayskip}{4pt plus 1pt minus 1pt}
\setlength{\belowdisplayskip}{4pt plus 1pt minus 1pt}
\setlength{\abovedisplayshortskip}{3pt}
\setlength{\belowdisplayshortskip}{3pt}
\newcommand{\step}[1]{\par\smallskip\noindent{\color{violetdeep}\bfseries #1}\hspace{0.35em}}
\newcommand{\ok}{\textcolor{violet}{\ensuremath{\checkmark}}}

% 紧凑节头：留白减半，并防止节头孤零零留在页底
\newcommand{\chdr}[2]{%
  \needspace{4\baselineskip}%
  \par\vspace{4pt}%
  \noindent\colorbox{violetdeep}{\color{white}\bfseries\small\;#1\;}\hspace{7pt}%
  {\color{violetdeep}\bfseries #2}\par\nobreak
  \vspace{1pt}{\color{petal}\hrule height 1pt}\vspace{4pt}}
\let\hdr\chdr

\begin{document}

\doctitle{热流从两端流入：边界条件为什么左负右正}
{§7.2 作业第 4 题（课件 p19）\ ·\ 杆长 $l$，两端各流入热流密度 $q_0$}
{依据课件 p19 题面、p14 力边界条件、p15 热流边界条件\quad 2026-11-09\quad 数学物理方法\,2}

\qstem{杆长 $l$，两端各有热流密度 $q_0$ 流入杆内（图上两个箭头都朝杆内）。请写出该问题的\textbf{边界条件}。\\[2pt]
{\small\color{inkgray}课件 p19 原文：\emph{The heat flux into the ends of the rod (whose length is $l$) is $q_0$.}}}

\hdr{§1}{认题：题目给的是热流密度，不是温度}

未知量 $u(x,t)$ 是杆内温度。课件 p13--p16 把边界条件分成三类；\textbf{第二类}（p14）给定的是边界上的法向导数 $\partial u/\partial n$，而不是函数值本身。本题给的正是端面的热流密度 $q_0$，量纲 $\mathrm{W/m^2}$，所以它落在第二类，且是非齐次的（所给的值不为零）。

\begin{center}
\begin{tabular}{c@{}c@{}c}
$q_0\;\longrightarrow$ & {\setlength{\fboxsep}{0pt}\fcolorbox{petal}{petal!45}{\rule{0pt}{0.45cm}\hspace{3.6cm}}} & $\longleftarrow\;q_0$\\[1pt]
 & {\small\color{inkgray}\makebox[3.6cm][s]{$x=0$\hfill$x=l$}} &
\end{tabular}\\[2pt]
{\small\color{inkgray}（杆，长 $l$，温度 $u(x,t)$；两侧箭头表示热流流入）}
\end{center}

图上看，两端都是“流入”，但两端的\textbf{外法向方向相反}。同一个“流入”，在两个端面上对应着相反的空间方向——这是全题唯一容易翻车的地方，符号就出在这里。

\hdr{§2}{唯一的工具：外法向 $+$ Fourier 定律}

课件 p15 的约定是取\textbf{外法向} $n$（从杆指向杆外），热流密度沿外法向的分量由 Fourier 定律给出：

\[
q_n=-k\,\frac{\partial u}{\partial n}
\]

其中 $k$ 是热导率（$\mathrm{W\cdot m^{-1}\cdot K^{-1}}$）。这句话给出两条操作规则：

\begin{itemize}
\item \textbf{“流入”意味着热流逆着外法向}，所以 $q_n<0$；量值为 $q_0$ 时就是 $q_n=-q_0$。这一条在两端都成立，因为“流入”参照的始终是杆的内部。
\item 把 $q_n$ 换回温度导数：$\dfrac{\partial u}{\partial n}=-\dfrac{q_n}{k}$，于是两端都有 $\dfrac{\partial u}{\partial n}=\dfrac{q_0}{k}$。
\end{itemize}

p15 已经替我们验过一遍右端：热流\textbf{流出}时 $\partial_xu|_{x=l}=-f/k$，\textbf{流入}时 $\partial_xu|_{x=l}=+f/k$，$f=0$ 则 $\partial_xu|_{x=l}=0$（绝热）。

\hdr{§3}{两端分别落笔：先写 $\partial u/\partial n$，再换到 $\partial u/\partial x$}

\step{统一形式（两端一样）}按 §2 的第二条，两个端面都满足

\[
\left.\frac{\partial u}{\partial n}\right|_{x=0}=\frac{q_0}{k},
\qquad
\left.\frac{\partial u}{\partial n}\right|_{x=l}=\frac{q_0}{k}.
\]

这一步之所以两端同号，是因为“外法向”本身就是随端面变的。

\step{换到 $x$ 导数（符号在这里分开）}左端 $x=0$ 的外法向指向 $-x$，所以 $\dfrac{\partial u}{\partial n}=-\dfrac{\partial u}{\partial x}$；代入上式：

\[
-\left.\frac{\partial u}{\partial x}\right|_{x=0}=\frac{q_0}{k}
\quad\Longrightarrow\quad
\left.\frac{\partial u}{\partial x}\right|_{x=0}=-\frac{q_0}{k}.
\]

右端 $x=l$ 的外法向指向 $+x$，所以 $\dfrac{\partial u}{\partial n}=+\dfrac{\partial u}{\partial x}$；代入同一条统一形式：

\[
\left.\frac{\partial u}{\partial x}\right|_{x=l}=\frac{q_0}{k}.
\]

也就是说，两端的差别不在“流入”这件事上，而在“外法向朝哪边”。

\hdr{§4}{答案}

\[
\boxed{\;\left.\frac{\partial u}{\partial x}\right|_{x=0}=-\frac{q_0}{k},
\qquad
\left.\frac{\partial u}{\partial x}\right|_{x=l}=+\frac{q_0}{k}\;}
\]

三种等价写法，考试写哪一种都算对，但同一份答案里只用一种：

\[
\text{外法向形式：}\ \left.\frac{\partial u}{\partial n}\right|_{\text{两端}}=\frac{q_0}{k};
\qquad
\text{通量形式：}\ -k\left.\frac{\partial u}{\partial x}\right|_{x=0}=q_0,\quad
-k\left.\frac{\partial u}{\partial x}\right|_{x=l}=-q_0 .
\]

通量形式最后一个是负号，读法是：在 $x=l$ 处，$+x$ 方向的净热流是 $-q_0$，也就是真实热流沿 $-x$ 流进杆内。

\hdr{§5}{四条自检}

\begin{itemize}
\item \textbf{量纲}：$\dfrac{q_0}{k}$ 的单位是 $\dfrac{\mathrm{W/m^2}}{\mathrm{W/(m\cdot K)}}=\mathrm{K/m}$，与 $\partial u/\partial x$ 一致 \ok{}。
\item \textbf{退化}：$q_0=0$ 时两端 $\partial_xu=0$，正是 p15 的绝热情形 \ok{}。
\item \textbf{温度形状}：两端都在吸热，端部比内部热，所以温度从左端往里降、从右端往里降，即左端 $u_x<0$、右端 $u_x>0$，与答案符号一致 \ok{}。
\item \textbf{能量}：对 $u_t=a^2u_{xx}$ 积分得 $\dfrac{\dd}{\dd t}\displaystyle\int_0^l \rho c\,u\dd x=k\big[u_x(l)-u_x(0)\big]=2q_0>0$，热量持续流入，温度整体上升 \ok{}。顺便说明：本题\textbf{没有稳态}——稳态要求 $u_{xx}=0$，即 $u$ 是直线、两端斜率相等，与 $-\frac{q_0}{k}\neq\frac{q_0}{k}$ 矛盾。
\end{itemize}

\hdr{§6}{和 p14 是同一副符号，不必背两套}

p14 的杆纵振动：两端受\textbf{向外}的力 $f$，Hooke 定律给 $f_n/S=E\,\partial u/\partial n$，于是两端 $\partial u/\partial n=f/(ES)$，换成 $x$ 导数也是“左端负、右端正”：

\[
\left.\frac{\partial u}{\partial x}\right|_{x=0}=-\frac{f}{ES},
\qquad
\left.\frac{\partial u}{\partial x}\right|_{x=l}=+\frac{f}{ES}.
\]

把 $f$ 换成 $q_0$、把 $ES$ 换成 $k$，就是本题。两题的共同结构是：\emph{同一个物理量作用在两个端面，写 $x$ 导数时符号必然相反，写外法向导数时符号相同。}

\hdr{§7}{写进作业本的三行}

\begin{enumerate}
\item 取外法向 $n$，写出 Fourier 定律 $q_n=-k\,\partial u/\partial n$。
\item 热流流入，$q_n=-q_0$；左端外法向沿 $-x$，右端沿 $+x$。
\item 代入得 $\left.\dfrac{\partial u}{\partial x}\right|_{x=0}=-\dfrac{q_0}{k}$，$\left.\dfrac{\partial u}{\partial x}\right|_{x=l}=+\dfrac{q_0}{k}$。
\end{enumerate}

\gapitem{错法一 · 两端写成同一个符号}{根源是把“流入”当成了同一个空间方向。判据是温度形状：两端吸热时温度是“两端高、中间低”，所以 $u_x$ 在左端为负、右端为正；写成两个 $-\frac{q_0}{k}$ 或两个 $+\frac{q_0}{k}$ 都与这个形状矛盾。}

\gapitem{错法二 · 系数写错}{把 $k$ 写成 $ES$ 是把热传导照抄成杆的纵振动。$ES$ 是拉压刚度，$k$ 是热导率；题目给的物理量换了，系数必须跟着换。漏写 $k$（写成 $\partial_xu=\mp q_0$）则是量纲不成立。}

\gapitem{错法三 · 用了内法向却没改符号}{若取内法向，则“流入”对应 $q_n=+q_0$，$\partial u/\partial n=-\frac{q_0}{k}$；结论一样，但中间每一步的符号都得跟着翻。用哪种法向都行，\textbf{前提是整篇统一}。}

\callout{\textbf{一句话带走}：边界条件的符号只由一件事决定——“流入”与所选法向的方向关系。取外法向后，两端流入都写成 $q_n=-q_0$；翻成 $\partial u/\partial x$ 时，左端才多出一个负号。}

\vspace{2pt}
{\small\color{inkgray}附注：本题只问边界条件。若要凑成完整的定解问题，还需一个初始条件 $u(x,0)=\varphi(x)$（温度分布），但题目没给，答案里不必臆造；热传导方程对时间是一阶的，初值只此一个。}

\end{document}
